\subsection{All-disjoint retained centers}
Use two factors $h=24$. For triple inputs write $G=\sum_TX_T$,
$T_i=\sum_{T\ni i}X_T$, and retain all $A_i=G-T_i$.
Each input triple avoids exactly 21 points, so
\[
 \sum_{i=1}^{24}A_i=21G,\qquad G=\frac1{21}\sum_i A_i.
\]
For a target triple $S$ the central correction is
\[
 G-\frac12\sum_{i\in S}A_i
       =\frac{\sum_{i\in S}T_i-G}{2}.
\]
An input triple meeting $S$ in $j$ points has central coefficient
$(j-1)/2$. Adding half the disjoint sum and subtracting half the
intersection-two sum leaves coefficient one exactly when $j=3$.
The scalar source-to-scatter transfer is therefore the identity.
The transparent old-value subtraction and new-value addition use this
identity only on fresh contributions and restore arbitrary old values.

The base-two paired exclusion and aligned pair-star producer retains
nested binary labels: common-pair supports have their orthonormal triple
span, other ordinary sums have their coordinate cover. Each retained
$A_i$ has the nondegenerate coordinate hyperplane excluding $i$,
of dimension 23. Hence all copied-center arguments and general
nondegenerate binary residuals of the preceding proof apply. The
selected local counts are
\[
 v=2024,\quad c_h=45764,\quad q_h=8120,\quad
 |\mathcal M|=8966,\quad R=44918,\quad\ell=552.
\]
The old internal histogram, starting at rank zero, is
\begin{small}\begin{align*}
[&8930,49042,30844,15972,11883,8272,6399,4356,6682,2263,\\
 &5109,1808,5317,1733,5391,1757,4908,3290,6233,4363,\\
 &4056,264,564,24,24].
\end{align*}\end{small}
It has rank mass $24R+2\ell$. Replace 24 rank-24 original transitions
by 24 rank-one transitions; all rank-23 copied transforms remain.
The new mass is $24R+\ell$.

\subsection{Physical endpoints and strict complex moment}
The copied-center two-stage complex contract in
\path{notes/copied-centers-complex.tex} is unchanged by the scalar
center decoding. For the tensor triple line $U$, let
$T=C_U$, $F=C_I$, $E=C_{U^\perp}=FT^{-1}$.
The two data outputs are $A=-Fy$ and $B=FT^{-2}x+Ey$.
Add $T^{-1}$ applied to a copy of $A$ into $B$, paying one rank-one
child per data pair. The inherited input/output $Z$ phases and
$i^{9f}$ scalar normalize $FT^{-2}$ to $F$; sign and bank corrections
complete the exact tensor $C$. Auxiliary source/sink operators remain
$C_{D_0}$ and $FC_{D_0}$ in the second stage. Every original role
therefore returns the exact selected-axis tensor operation on arbitrary
inputs. This algebraic statement holds over $\mathbb Q(i)$ as well as
over Gaussian dyadics.

With $m_c=576$, $N_c=v^2=4096576$, $B_c=vR=90914032$, put
\[
 W_c=2N_c+2B_c=190021216,\quad L_c=2v\ell=2234496,
\]
\[
 s_c=W_cm_c-N_c+L_c=109450358336,\qquad N_c-L_c=1862080.
\]
The complete list is $2B_c$ width-552 auxiliary exteriors,
$2N_c$ width-529 data macros, $2vH'_r$ width-$r$ internal children,
$4N_c$ width-23 data fronts, and $N_c$ width-one endpoint copies.
These disjoint classes include both kinds of paid copies and sum to
$s_c$. All binary residuals, including alternating forms, use the retained
one-child Gauss normal form with paid bit adapters and unit phases.
At $a_c=1949/25000000$ the exact rational enclosure gives
\[
 \sum_r\frac{n_r}{W_c}(r/m_c)^{1-a_c}
 <1-1.3338373969\times10^{-10}.
\]

\subsection{One exact odd-denominator grid}
The only new odd divisor is 21. Let $d$ bound the root axis count,
$D$ bound its maximum complex recursion depth, and $G_0$ bound the
number of grouped scalar operations per node. Choose
\[
 K=G_0(D+1),\qquad \text{common grid }2^{-P}21^{-K}.
\]
Root inputs and arbitrary dirty auxiliaries start Gaussian dyadic.
Multiply their integer numerators by $21^K$ once at the outer boundary.
Every local coefficient contains at most one factor 21 in its denominator,
so a local finite program adds at most $G_0$ to the odd-denominator
exponent. A completed $u$-axis child has Gaussian-dyadic coefficients
and therefore preserves its input's existing odd-denominator exponent;
it adds no new odd factor. Only the active chain of unfinished calls
can accumulate new factors. Its length is at most $D+1$, proving that
all intermediates lie on the common grid and every local division by
21 is an exact integer division of its encoded numerator.

This divisibility claim uses the root dyadic-input invariant, not an
assumption about arbitrary numerators on the finest grid. Internal values
may have odd denominators. They are neither rounded nor re-encoded when
a child returns. At the completed outer return the operator is the exact
Gaussian-dyadic tensor on root dyadic inputs, so the final numerators
are divisible by $21^K$; exact unscaling restores dyadic encoding.

Since the largest child is $552<576$, $D=O(\log d)$ and
$K=O(\log d)$ with fixed constants. Constant-divisor long division
is linear in numerator length. Initial and final scaling by $K$
successive constant multiplications/divisions cost $O(V\log d)$ and
occur only at the outer boundary. The extra numerator length is
$O(\log d)$, preserving the retained $O(p)$ word bound. This overhead
is subordinate to the positive-power complex layer bound.

For explicit semantic constants use
\[
 g_h=4(c_h+v)+10v+4hv+4h^2+8h+8=408200,
\]
\[
 G_0=N_c+2v(g_h+2h)=1656684480,
\]
where endpoint additions and copied-center scalar work are included.
The complete child contributes at most $u$ binary denominator bits
and row-norm factor $2^u$, with no additional odd exponent. A local
21-division contributes fewer than five denominator-height bits,
covered by the fixed scalar charge. Set
\[
\begin{aligned}
 E_0&=64(W_c+m_c+G_0+1)^3=403063477262941121443761178688,\\
 \mathcal E&=2G_0W_c^2+8s_c+4W_c+4+32m_c\\
           &=119639333593035096270287748<E_0,\\
 B_0&=s_c+E_0,\qquad C_0=32m_cB_0^2.
\end{aligned}
\]
With $f=\lfloor e/m_c\rfloor$, the completed-child induction is
\[
 A(e)\le A(552f)+s_cf+E_0\le2B_0e,
\]
since $2B_0(m_c-552)\ge s_c+E_0$. It counts total denominator height
and magnitude, including the local odd factor. The retained prefix and
normalization allowance $18e$ fits inside $C_0e$. Together with the
$O(\log d)$ common-grid encoding overhead this gives $C_1=1$ and
$O(p)$ records for $d=\lfloor b_{\rm in}^{\epsilon}\rfloor$,
$p=6b_{\rm in}$, fixed $\epsilon<1$.
